\section{An exact energy reduction under a covariance certificate}

Consider a nonempty linearly separable labeled set, and put
$z_i=y_i\bar x_i$. Let $a$ be its minimum-norm affine unit-margin
separator and $Q=\|a\|^2$. Suppose there are probabilities $\rho_i$ on
active constraints such that
\begin{equation}\label{eq:certificate}
 z_i\cdot a=1\quad(\rho_i>0),\qquad
 \sum_i\rho_i z_i=a/Q,\qquad
 C:=\sum_i\rho_i z_i z_i^T\preceq2I.
\end{equation}
The first two conditions are the usual normalized optimality
certificate. The covariance condition is an additional assumption;
it will be verified explicitly for the histories used below.
It also implies $Ca=a/Q$.

Write $b_j=Nv_j$. The energy and output are
\begin{equation}\label{eq:energy}
 E_N=\frac1N\sum_j(d\|u_j\|^2+b_j^2),\qquad
 F(x)=\frac1N\sum_j b_j\phi_\epsilon(u_j\cdot\bar x).
\end{equation}

\begin{theorem}[Exact energy and equality]\label{thm:exact}
Assume \eqref{eq:certificate} and
$d\ge4\max\{1,\epsilon^{-6}\}$. At every hidden width $N\ge1$, the
minimum energy over networks satisfying the training margins
$y_iF(x_i)\ge1/2$ is
\begin{equation}\label{eq:profile}
 J_d(Q)=\min_{r\ge0}
 \left\{dQr^2+\frac1{4(\epsilon+\tanh r)^2}\right\}.
\end{equation}
The unique minimizing radius is positive and satisfies
\begin{equation}\label{eq:stationarity}
 4dQr(\epsilon+\tanh r)^3=\operatorname{sech}^2r.
\end{equation}
The profile is strictly increasing in $Q$ and is below
$1/(4\epsilon^2)$. It is attained by identical neurons
\[
 u_j=ra,\qquad b_j=\frac1{2(\epsilon+\tanh r)}.
\]
At equality every neuron has this form or has both signs reversed.
Consequently the entire output function is unique. The same assertions
hold for arbitrary probability distributions over neurons in place of
the uniform distribution on $N$ neurons.
\end{theorem}

\begin{proof}
Put $h=\epsilon+\tanh r$. The scalar objective has derivative
$2dQr-\operatorname{sech}^2r/(2h^3)$, which is strictly increasing
from a negative value to a positive value. This proves the existence
and uniqueness of the positive optimum and \eqref{eq:stationarity}.
Define
\[
 \beta=\operatorname{sech}^2r,\qquad
 \alpha=\epsilon+\tanh r-r\beta,\qquad
 \lambda=h^{-2},\qquad K=\lambda^2/4,\qquad M=1+\epsilon.
\]
Then $0<\alpha\le M$, $0<\beta\le1$,
$h=\alpha+\beta r$, and $K\le1/(4\epsilon^4)$.
The inequality $\alpha\ge\epsilon$ follows by differentiating
$\tanh r-r\operatorname{sech}^2r$, whose derivative is nonnegative.
Stationarity gives $dQr=K\beta h$, so
\begin{equation}\label{eq:quadratic-floor}
 dQ t^2-K(\alpha+\beta t)^2\ge
 dQr^2-Kh^2=:\zeta=-\frac{\alpha}{4h^3}<0
 \quad(t\ge0).
\end{equation}
Indeed its leading coefficient is
$dQ-K\beta^2=K\alpha\beta/r>0$ and its minimum is at $t=r$.

For a neuron $u$, decompose $u=ta+w$ with $w\perp a$, and regard
$s=z_i\cdot u$ as a random variable under $\rho$.
The active identities imply
\[
 \E_\rho s=t,\qquad
 \operatorname{Var}_\rho(s)=w^TCw\le2\|w\|^2.
\]
Since $\sup_s|\tanh''s|<1$, Taylor's theorem gives
\[
 |\E_\rho\tanh s-\tanh t|
 \le\tfrac12\operatorname{Var}_\rho(s)\le\|w\|^2.
\]
For $S(u)=\E_\rho\phi_\epsilon(s)$, the threshold part has
absolute average at most $\epsilon$. Concavity of tanh on the
positive half-line therefore yields
\begin{equation}\label{eq:atom-average}
 |S(u)|\le\epsilon+|\tanh t|+\|w\|^2
       \le\alpha+\beta|t|+\|w\|^2.
\end{equation}
Also $|S(u)|\le M$.

For any readout $b$, complete the square:
\begin{equation}\label{eq:square}
 d\|u\|^2+b^2-\lambda bS(u)
 \ge d\|u\|^2-KS(u)^2.
\end{equation}
If $\|u\|^2\ge KM^2/d$, the right side is nonnegative and hence
strictly greater than $\zeta$. Otherwise set $q=\|w\|^2$ and
$R_0=\sqrt K M/\sqrt d$. Then $q\le R_0^2$ and
$|t|\le R_0/\sqrt Q$. By \eqref{eq:atom-average}, the right side of
\eqref{eq:square} is at least
\begin{equation}\label{eq:orthogonal-cost}
 \bigl[dQt^2-K(\alpha+\beta|t|)^2\bigr]
 +q\bigl[d-2K(\alpha+\beta|t|)-Kq\bigr].
\end{equation}
Any unit-margin constraint gives $Q\ge1/(d+1)$, so $dQ\ge1/2$.
Consequently
\[
 2K(\alpha+\beta|t|)\le2KM(1+\sqrt{2K}).
\]
If $0<\epsilon\le1$, this is at most
$\epsilon^{-4}(1+1/(\sqrt2\epsilon^2))\le2\epsilon^{-6}\le d/2$.
If $\epsilon\ge1$, use $M\le2\epsilon$ to bound it by
$\epsilon^{-3}(1+1/(\sqrt2\epsilon^2))\le2\le d/2$.
Similarly,
\[
 Kq\le K^2M^2/d\le
 \begin{cases}
 1/(4d\epsilon^8)\le d/64,&0<\epsilon\le1,\\
 1/(4d\epsilon^6)\le d/64,&\epsilon\ge1.
 \end{cases}
\]
Thus the second bracket in \eqref{eq:orthogonal-cost} is at least
$31d/64>0$. Equations \eqref{eq:quadratic-floor}--\eqref{eq:orthogonal-cost}
prove the atomwise floor $\zeta$ for every $u,b$.

Averaging over neurons and using the signed training margins gives
\[
 E_N\ge\lambda\E_j[b_jS(u_j)]+\zeta
      \ge\lambda/2+\zeta
      =dQr^2+\frac1{4h^2}.
\]
The cloned construction attains equality: every signed training
preactivation is at least $r$, and active ones equal $r$.
The same proof integrates over any neuron probability law with finite
energy.

Equality forces $w=0$, $|t|=r$, and equality in square completion.
Thus $u=\pm ra$ and $b=\pm1/(2h)$ with matching signs.
Oddness of the activation makes all such atoms have the same output.
Finally, the negative derivative at zero puts the optimum below
$1/(4\epsilon^2)$. For $Q_1>Q_0$, evaluation at the positive
minimizer $r_1$ for $Q_1$ gives
$J_d(Q_1)-J_d(Q_0)\ge d(Q_1-Q_0)r_1^2>0$.
\end{proof}

\section{A positive energy gap uniform over hidden width}

The next lemma addresses a subtlety of the online order: only the first
new label of a block has been revealed when its error is counted.
That intermediate history need not have the covariance certificate.

\begin{lemma}[Uniform gap to a different label]\label{lem:uniform-gap}
Suppose a history satisfies Theorem~\ref{thm:exact}, its affine
minimizer $a$ has nonzero scores on the entire cube, and its scalar
output predicts the wrong label at a query $x_0$. If some globally
margin-realizable function fits the history and the correct query
label, there exists $e_1>J_d(Q)$ such that every network fitting those
constraints and global margins has energy at least $e_1$, at every
width.
\end{lemma}

\begin{proof}
We use a closed relaxation that includes all finite-width networks and
all limits of their neuron distributions. A relaxed neuron is
$(u,b,H)$, where $H\in[-1,1]^X$ and
\[
 H(x)=\sgn(u\cdot\bar x)\quad\text{whenever }u\cdot\bar x\ne0.
\]
This is a closed subset of a finite-dimensional Euclidean space.
Its output is the continuous function of the relaxed coordinates
\[
 b\bigl[\tanh(u\cdot\bar x)+\epsilon H(x)\bigr].
\]
At zero preactivation the relaxation permits any jump value in
$[-\epsilon,\epsilon]$. The proof of Theorem~\ref{thm:exact} still
applies: its only use of the jump term in \eqref{eq:atom-average} was
the bound on its absolute average by $\epsilon$.

Consider all probability measures on relaxed neurons satisfying the
closed training margins and the correct signed query margin at least
$1/2$. We may omit global constraints at other inputs, which only
enlarges this class. It is nonempty by hypothesis. The infimum of
mean energy $\E[d\|u\|^2+b^2]$ is attained. To see this, take a
minimizing sequence with a common finite energy bound. The second
moments make it tight, while the $H$ coordinates lie in a compact set.
A weakly convergent subsequence is supported on the same closed
relaxed-neuron space. The energy is lower semicontinuous. Each output
constraint passes to the limit: its integrand has magnitude at most
$M|b|$, and the second-moment bound makes these integrands uniformly
integrable. Truncating $b$ gives bounded continuous test functions,
which proves the required convergence of means.

Theorem~\ref{thm:exact} puts the attained energy at least $J_d(Q)$.
If equality held, its atomwise equality proof would force almost every
neuron to have $u=\pm ra$ and the matching readout. These rows have no
zero preactivation anywhere on the cube. All their jump values are
therefore forced, and every neuron has the same scalar output, which
is wrong at $x_0$. This contradicts the correct query margin. The
attained energy is consequently some $e_1>J_d(Q)$. Every finite-width
network is a uniform empirical measure in this relaxation, proving the
width-uniform claim.
\end{proof}
